\chapter*{Preface}
\addcontentsline{toc}{chapter}{Preface}

\section*{A different promise}

Most calculus textbooks make you a quiet promise: memorize these formulas, practice these
templates, and you will pass. This book takes a different path. Work through these ideas with pencil
in hand, and you will \emph{understand}. You will see why the derivative takes the form it does,
where it came from, what a limit actually asserts, why the chain rule holds together geometrically,
and why the area under a curve and the slope of a curve turn out to be two sides of a single fact.
Beyond calculating answers, you will look at a symbol like $\dd x$, or a phrase like ``the limit as
$x$ approaches $a$,'' and see what is actually happening underneath.

That is a demanding promise to keep, and it asks more of you. Understanding cannot be poured into an
unresisting mind; it has to be constructed from direct experience. So this book is built, brick by
brick, from the ground up. We start with what a function truly represents, and we conclude, many
chapters later, with differential equations -- the language in which physics, biology, and economics
write down how the world evolves. Every step in between is crafted to feel natural, as though you
could have discovered it yourself while exploring a concrete problem.

\begin{focusbox}
The single sentence that organizes this entire book: \emph{calculus is the mathematics of change and
accumulation, and the bridge between them is the limit.} If you keep that sentence in view, every
chapter will find its natural place.
\end{focusbox}

\begin{figure}[htbp]
\centering
\begin{tikzpicture}
  % Panel (a): The Derivative (Rate of change at an instant)
  \begin{scope}[shift={(0,0)}]
    \node[anchor=south west, font=\scriptsize\scshape, text=sourceblue] at (0,4.55) {Panel (a): The Rate Problem (Derivative)};
    \draw[->, thick, -stealth] (-0.3,0) -- (4.8,0) node[right, font=\footnotesize] {$x$};
    \draw[->, thick, -stealth] (0,-0.3) -- (0,4.2) node[above, font=\footnotesize] {$y$};
    
    % Curve y = 0.2*x^2 + 0.5
    \draw[curveblue, very thick, domain=0.4:4.2, smooth] plot (\x, {0.2*(\x)^2 + 0.5});
    \node[curveblue, font=\small, right] at (3.9, 3.8) {$y = f(x)$};
    
    % Point P at x = 2
    \draw[tangentred, thick, domain=0.6:3.8] plot (\x, {0.8*\x - 0.3});
    
    % Tangent triangle (nudges dx and df)
    \draw[dashed, gray] (2, 1.3) -- (3.2, 1.3) node[midway, below, font=\footnotesize, text=black] {$\dd x$};
    \draw[dashed, gray] (3.2, 1.3) -- (3.2, 2.26) node[midway, right, font=\footnotesize, text=black] {$\dd f$};
    \filldraw[sourceblue] (2, 1.3) circle (2pt) node[above left, font=\footnotesize, text=black] {$(x, f(x))$};
    
    % Slope callout box
    \node[draw=tangentred!70, fill=tangentred!10, rounded corners=2pt, font=\footnotesize, align=center] at (1.8, 3.3)
      {Instantaneous slope:\\[0.2em] $\dfrac{\dd f}{\dd x} = \lim_{\Delta x \to 0} \dfrac{\Delta f}{\Delta x}$};
  \end{scope}

  % Divider corridor
  \draw[gray!30, dashed, thick] (5.7,-0.2) -- (5.7,4.7);

  % Panel (b): The Integral (Accumulation across an interval)
  \begin{scope}[shift={(6.4,0)}]
    \node[anchor=south west, font=\scriptsize\scshape, text=targetgreen] at (0,4.55) {Panel (b): The Accumulation Problem (Integral)};
    \draw[->, thick, -stealth] (-0.3,0) -- (4.8,0) node[right, font=\footnotesize] {$x$};
    \draw[->, thick, -stealth] (0,-0.3) -- (0,4.2) node[above, font=\footnotesize] {$y$};
    
    % Shaded area under curve from x = 1 to x = 3.8
    \fill[areafill!35] (1,0) -- plot[domain=1:3.8, smooth] (\x, {0.2*(\x)^2 + 0.5}) -- (3.8,0) -- cycle;
    
    % Curve y = 0.2*x^2 + 0.5
    \draw[curveblue, very thick, domain=0.4:4.2, smooth] plot (\x, {0.2*(\x)^2 + 0.5});
    \node[curveblue, font=\small, right] at (3.9, 3.8) {$y = f(x)$};
    
    % One representative rectangle slice at x = 2.2 to 2.6
    \filldraw[fill=quotientorange!25, draw=quotientorange!80, thick] (2.2, 0) rectangle (2.6, 1.65);
    \draw[<->, >=stealth, gray!80] (2.2, -0.2) -- (2.6, -0.2) node[midway, below, font=\footnotesize, text=black] {$\dd x$};
    
    % Area slice callout
    \node[draw=quotientorange!80, fill=quotientorange!10, rounded corners=2pt, font=\footnotesize, align=center] at (1.8, 3.3)
      {Accumulated sum:\\[0.2em] $\mathrm{Area} = \int_a^b f(x)\,\dd x$};
    
    \node[font=\footnotesize] at (1, -0.25) {$a$};
    \node[font=\footnotesize] at (3.8, -0.25) {$b$};
  \end{scope}
\end{tikzpicture}
\caption{The two grand pillars of calculus. Left: The rate problem zooms into a single point to find the local steepness $\dd f/\dd x$. Right: The accumulation problem sums up infinitely many thin slices of area $f(x)\,\dd x$ across an interval. The central discovery of calculus is that differentiation and integration invert each other.}
\label{ch00:two-problems}
\end{figure}

There is a compelling reason to insist on understanding beyond intellectual satisfaction.
Techniques acquired only by rote memorization evaporate. Within a year of finishing a procedural
course, many students cannot integrate by parts or state the Mean Value Theorem, because nothing
anchored the rules in place -- they were an arbitrary list, and unanchored lists decay. But ideas
you have explored visually and conceptually do not vanish in that way, because each one is supported
by the geometric intuitions around it. If you forget the quotient rule, you can sketch a small
nudge and rederive it in two lines. If you forget the formula for a Taylor polynomial, you can
reconstruct the matching-derivative requirement that produced it. True understanding is never slower
than memorization; it is the only thing that endures.

\section*{How this book is built: intuition, then picture, then proof}

Every important idea in this book is introduced in the same three movements, and it will help you to
know the rhythm in advance.

\textbf{First, intuition.} Before any symbols, we ask what the idea \emph{is} in plain human terms.
What is a derivative? It is how fast something is changing right now -- the speedometer reading, not
the trip average. What is an integral? It is a total accumulated from infinitely many tiny
contributions -- the full distance built up from every instant's speed. We will state the idea the
way you might explain it to a sharp friend over coffee, with no apology and no formula.

\textbf{Second, the picture.} Calculus was discovered with pictures, and it is still understood with
pictures. A derivative is the slope of a tangent line. An integral is an area. A limit is a process
of closing in. Almost every idea in this book has an exact geometric meaning, and we will draw it.
You will see secant lines pivot down into a tangent; you will see rectangles thicken into a smooth
region of area; you will see the graphs of Taylor polynomials creep outward to hug a curve. The
figures are not decoration. They are arguments you can see.

\begin{figure}[htbp]
\centering
\begin{tikzpicture}
\begin{axis}[
    width=11cm, height=7cm,
    axis lines=middle,
    xlabel={$x$}, ylabel={$y$},
    xmin=-0.3, xmax=2.6, ymin=-0.3, ymax=4.6,
    xtick=\empty, ytick=\empty,
    domain=0:2.4, samples=100,
    clip=false,
]
% the curve y = x^2
\addplot[curveblue, very thick] {x^2};
\node[curveblue, font=\small] at (axis cs:2.15,4.2) {$y=x^2$};
% point of tangency at x=1
\addplot[only marks, mark=*, mark size=2pt, sourceblue] coordinates {(1,1)};
\node[font=\footnotesize, left=2pt] at (axis cs:1,1) {$(1,1)$};
% secant lines pivoting toward the tangent at x=1.
% A secant from (1,1) to (a,a^2) has slope a+1; here a = 2.2, 1.8, 1.4.
\addplot[gray!60, thin, domain=0.2:2.5] {3.2*(x-1)+1};
\addplot[gray!80, thin, domain=0.2:2.5] {2.8*(x-1)+1};
\addplot[gray, thick, domain=0.2:2.5] {2.4*(x-1)+1};
% second intersection points sliding in toward (1,1)
\addplot[only marks, mark=*, mark size=1.5pt, gray] coordinates {(2.2,4.84) (1.8,3.24) (1.4,1.96)};
% the tangent (slope 2) -- the limiting line
\addplot[tangentred, very thick, domain=-0.1:2.5] {2*(x-1)+1};
\node[tangentred, font=\small, right] at (axis cs:2.35,3.6) {tangent line};
\node[gray!70!black, font=\small, left] at (axis cs:0.35, -0.2) {secants pivot down};
\end{axis}
\end{tikzpicture}
\caption{The core geometric intuition of the derivative. Each gray secant line cuts the curve at two points; as the second point slides toward the point of tangency, the secants pivot and close in on a single red line. That limiting line is the tangent, and its slope is the derivative. We will build this picture carefully in the chapters on limits and the derivative.}
\label{ch00:secant-tangent}
\end{figure}

\textbf{Third, the proof -- clear and complete, but never dry.} Intuition can mislead, and pictures can
hide special cases. So after we see why an idea should be true, we make it precise and we prove it.
But a proof in this book is never a wall of symbols dropped on you from a height. We always explain
the \emph{idea} of the proof first -- the one move that makes it work -- and then we walk through the
steps narrating as we go, the way a good teacher talks through a derivation at the board. A proof is
just a careful retelling of why something has to be true. If you cannot feel why each line follows,
the proof has failed you, and we will try not to fail you.

\begin{intuitionbox}
You can read this book at two depths. On a first pass, read the intuition and study the pictures,
and let the proofs wash over you to get the shape. On a second pass, slow down and follow every line
of the derivations. Both passes are legitimate. Mathematics is rarely understood in one sitting, and
nobody -- not Newton, not you -- absorbs it linearly. Circling back is not failure; it is the method.
\end{intuitionbox}

\section*{The colored boxes: a visual language}

Throughout the book you will meet recurring colored boxes. Each one plays a fixed role, so once you
learn the color, you know what kind of thinking is coming. Three of them carry the special spirit of
this book, and they deserve introduction by name.

\begin{historybox}
A \textbf{History} box tells you where an idea came from, and -- just as important -- what its
inventors got right by instinct before anyone could justify it. Calculus was used brilliantly for a
century and a half before it was made rigorous. Newton and Leibniz computed orbits and areas with
``infinitesimals'' -- quantities supposed to be smaller than any positive number yet not quite zero
-- and they were savagely, correctly criticized for it (the philosopher George Berkeley mocked these
vanishing quantities as ``the ghosts of departed quantities''). And yet the answers were right.
History boxes show you this drama: the right intuition arriving first, the rigorous foundation
arriving later, with Cauchy and Weierstrass finally pinning down the limit in the 1800s. Knowing the
story tells you which ideas were hard-won, and reassures you that your own confusion is in good
company.
\end{historybox}

\begin{happeningbox}
A \textbf{What is really happening} box pulls back the curtain on a symbol or a computation. What
\emph{is} that $\dd x$ in $\int f(x)\,\dd x$ -- a number, a ghost, a leftover? What does
$\lim_{x\to a}$ actually claim, when nothing ever reaches $a$? What does $e^x$ truly mean, and why
does it equal its own derivative? These boxes exist because notation is compressed thought, and the
compression can hide the very thing you most need to see. We open the box and look inside.
\end{happeningbox}

\begin{whybox}
A \textbf{Why the technique works} box explains the underlying mechanism behind a method rather than merely cataloging steps.
Every computational rule in calculus is a disguised truth about limits or about functions, and the
disguise is usually thin once you know to look. Substitution, for example, is nothing more than the
chain rule run backwards; integration by parts is the product rule run backwards. When you see
\emph{why} a technique works, it stops being a trick you have to remember and becomes a consequence
you could have predicted. That is the difference between a magician and a physicist watching the
same card trick.
\end{whybox}

The remaining boxes round out the conversation, and you will get to know them by sight:

\begin{itemize}
  \item An \textbf{Intuition} box (green) states an idea in plain language before we make it precise.
  \item A \textbf{Common trap} box (red) flags a mistake almost everyone makes -- a place where the
  intuition lies, or where the notation invites an error. Forewarned is forearmed.
  \item A \textbf{Walking question} box (orange) hands you a question to chew on away from the page,
  on a walk or in the shower. The best calculus thinking often happens when you are not staring at
  symbols. These have no answers printed; they are yours to turn over.
  \item A \textbf{Work it yourself} box (purple) asks you to compute one small thing by hand, right
  then, because some truths only land in the fingers. Reading about a derivative is not the same as
  taking one.
  \item An \textbf{Application} box (magenta) shows the idea doing real work in physics, biology,
  economics, or machine learning -- where the abstraction earns its keep.
  \item A \textbf{Connection} box (teal) links the idea to other parts of mathematics, so the
  subject feels like one landscape rather than a list of topics.
  \item A \textbf{Why this matters} box (cyan) steps back to say why an idea is beautiful or
  central -- the aesthetic and structural payoff.
  \item A \textbf{Focus} box (blue) marks the one thing in a section you should not lose.
\end{itemize}

\begin{marginnotebox}
You do not have to read the boxes in any particular order relative to the main text, and you will
not break anything by skipping one and coming back. They are asides in the original sense: voices
leaning in from the margin. The main argument always stands on its own in the running prose.
\end{marginnotebox}

\section*{The arc of the book: from functions to differential equations}

It helps to see the whole journey before you take the first step, the way you might glance at a
trail map before a long hike. Here is the territory, and why it is laid out the way it is.

\textbf{We begin with functions and the objects calculus acts on.} Calculus does not study numbers;
it studies \emph{functions} -- rules that turn an input into an output, machines for relating one
changing quantity to another. So before anything else we make sure you truly understand functions,
and especially the two families that calculus cannot live without: the exponential functions
(and the remarkable number $e$, which we will meet from three different directions -- compound
interest, a limit, and the function that is its own rate of change) and the trigonometric functions
(which, properly understood, measure rotation and periodic oscillation alongside static geometry). These early
chapters are about \emph{meaning}: what these objects are, before we ever differentiate them.

\textbf{Then we build the limit, the keystone of everything.} The limit is the idea that makes the
whole subject possible and the idea that took mathematicians the longest to get right. We approach
it first through pure intuition and a long historical look at the centuries of productive confusion,
and only then through the precise ``$\eps$--$\delta$'' definition that finally made it rigorous. We
spend real time here because everything downstream -- derivatives, integrals, series, all of it --
is secretly a statement about limits. Rush the limit and the rest stays foggy; understand it and the
rest comes into focus.

\textbf{With limits in hand, we build the derivative.} The derivative is the limit of those pivoting
secant lines in Figure~\ref{ch00:secant-tangent}: instantaneous rate of change, the slope of the
tangent. We learn what it means, then we earn the rules -- the power rule, product and quotient
rules, and the chain rule -- by deriving each from the definition, so they are never arbitrary. Then
we put the derivative to work: linear approximation, the Mean Value Theorem (the quiet theorem that
secretly justifies half of calculus), finding maxima and minima, optimization, and l'H\^opital's
rule.

\textbf{Next we build the integral, and meet the great theorem.} An integral is an accumulation -- an
area, a total, a sum of infinitely many infinitesimal pieces. We define it carefully as a limit of
sums, and then we arrive at the summit of the first half of the book: the Fundamental Theorem of
Calculus, which reveals that differentiation and integration are inverse operations. The slope
problem and the area problem, which look completely unrelated, turn out to be the same problem read
forwards and backwards. From there we develop the techniques of integration -- substitution, parts,
partial fractions -- always with a ``why it works'' box exposing the rule underneath the trick.

\textbf{Then we extend the reach: series, and then more than one variable.} We learn to add up
infinitely many terms (sequences and series), and we discover Taylor series, which let us rebuild
familiar functions like $e^x$, $\sin x$, and $\cos x$ out of nothing but their derivatives at a
single point -- one of the most beautiful ideas in all of mathematics. Then we step up a dimension:
functions of several variables, partial derivatives, the gradient, multiple integrals, and a first
taste of vector calculus.

\textbf{Finally we arrive at differential equations.} This is where calculus turns back toward the
world it came from. A differential equation is a sentence about how a quantity changes, written in
the language of derivatives, and its solution is the quantity's whole life story. Population growth,
radioactive decay, a swinging pendulum, a cooling cup of coffee, an electrical circuit -- all of
them are differential equations. Reaching them is the point of the climb.

\begin{figure}[h]
\centering
\begin{tikzpicture}[
    node distance=6mm,
    stage/.style={draw, rounded corners, align=center, font=\small,
                  minimum height=8mm, inner sep=3pt},
    arr/.style={-{Stealth[length=2.5mm]}, thick, gray}
]
\node[stage, fill=sourceblue!15, draw=sourceblue] (fn) {Functions\\ $e^x$, $\sin$, $\cos$};
\node[stage, fill=structurepurple!15, draw=structurepurple, right=of fn] (lim) {Limits};
\node[stage, fill=tangentred!12, draw=tangentred, right=of lim] (der) {Derivatives};
\node[stage, fill=areafill!35, draw=curveblue, below=8mm of der] (int) {Integrals \&\\ the FTC};
\node[stage, fill=quotientorange!18, draw=quotientorange, left=of int] (ser) {Series \&\\ Taylor};
\node[stage, fill=targetgreen!18, draw=targetgreen, left=of ser] (multi) {Several\\ variables};
\node[stage, fill=structurepurple!18, draw=structurepurple, below=8mm of multi] (ode) {Differential\\ equations};
\draw[arr] (fn) -- (lim);
\draw[arr] (lim) -- (der);
\draw[arr] (der) -- (int);
\draw[arr] (int) -- (ser);
\draw[arr] (ser) -- (multi);
\draw[arr] (multi) -- (ode);
\end{tikzpicture}
\caption{The arc of the book. Each idea rests on the one before it. The limit is the keystone: every
arrow after it is, underneath, a statement about limits. The destination -- differential equations --
is where calculus describes how the world actually changes.}
\label{ch00:arc}
\end{figure}

\begin{aestheticbox}
Notice the shape of Figure~\ref{ch00:arc}. It is not a list of independent topics to be ticked off.
It is a single structure where each stone holds up the next, and the limit is the keystone in the
arch. This is why we refuse to rush the early chapters. A keystone set carelessly brings down
everything resting on it.
\end{aestheticbox}

\section*{How to read this book}

A few practical suggestions, earned from watching many people learn this subject.

\textbf{Draw.} This is the single most important piece of advice in the preface, and it is why it
gets its own paragraph. Keep paper and pencil beside you and redraw the figures yourself -- not trace
them, draw them. Sketch the secant lines pivoting into the tangent. Shade the area under a curve and
chop it into rectangles. Plot a function and guess its derivative before you compute it. The act of
drawing forces you to decide what you actually believe about the shape, and it exposes vague
understanding instantly. Mathematicians who look like they are ``just thinking'' are very often
drawing in their heads, and they learned to do that by drawing on paper first.

\begin{walkbox}
Here is your first walking question, to carry away from this preface. A car's speedometer reads
$60$ miles per hour at a single instant -- say, exactly at noon. But speed is distance divided by
time, and at a single \emph{instant} no time passes and no distance is covered: the car is, in that
frozen instant, nowhere going anywhere. So what could ``$60$ miles per hour \emph{right now}''
possibly mean? Sit with that. The real resolution of this paradox is the derivative, and it is
also, in miniature, the central idea of the whole first half of this book.
\end{walkbox}

\textbf{Ask ``why'' relentlessly, and refuse to be satisfied too early.} When the book states a rule,
stop and ask why it should be true before you read on. When it draws a picture, ask what would change
if the curve bent the other way. When a proof makes a move, ask why that move and not another. This
book is written to reward that interrogation -- nearly every ``why'' you can think to ask is answered
somewhere, often in a colored box -- but the answer means far more if you have felt the question
first. A question you have actually wrestled with is a hook the answer can hang on. An answer to a
question you never asked just slides off.

\textbf{Compute things by hand, especially when you would rather not.} There is a kind of
understanding that only arrives through your fingers. When a ``Work it yourself'' box asks you to
take one derivative or evaluate one integral, do it then and there, even though you could read past
it. The friction is the point.

\begin{labbox}
Try this one now, before you turn the page, with nothing but arithmetic. Take the function that
squares its input, $f(x)=x^2$. Compute the average rate of change between $x=1$ and a nearby point
$x=1+h$ for shrinking $h$: that average rate is
\[
\frac{f(1+h)-f(1)}{h}=\frac{(1+h)^2-1}{h}=\frac{2h+h^2}{h}=2+h.
\]
Now plug in $h=0.1$, then $h=0.01$, then $h=0.001$. You get $2.1$, then $2.01$, then $2.001$. The
numbers are visibly homing in on $2$. That $2$ is the slope of the red tangent line in
Figure~\ref{ch00:secant-tangent}, and you just computed your first derivative -- by watching a quotient
close in on a value it never quite reaches. Hold on to the feeling of that ``closing in.'' It is the
limit, and the whole book turns on it.
\end{labbox}

\textbf{Do the exercises.} Every chapter ends with a set of them -- some computational, some
conceptual, and usually a proof or two. The computational ones build fluency; the conceptual ones
build understanding; the proofs build the ability to be \emph{sure}. You do not have to do all of
them, but do not do none of them. Mathematics is not a spectator sport, and reading a worked example
feels deceptively like doing one. It is not the same, in the way that watching someone swim is not
the same as getting wet.

\begin{trapbox}
The most common way to fail at calculus is to read it like a novel -- smoothly, continuously, never
stopping. It cannot be read that way. If you find yourself nodding along for three pages without
pausing, you have almost certainly stopped understanding and started skimming. Real reading of
mathematics is slow and full of stops: rereading a line, working a tiny case in the margin, drawing
the picture. If a sentence does not make sense, do not push past it hoping it will clear up. Stop.
Draw something. Try a number. The fog is information.
\end{trapbox}

\begin{recoverybox}
And when you get stuck -- you will get stuck, everyone gets stuck -- that is not a sign you are bad at
this. Getting stuck is what doing mathematics feels like from the inside, for everyone, at every
level. The difference between people who learn calculus and people who give up is almost never raw
talent. It is whether, on getting stuck, they draw a picture, try a smaller case, reread the
definition, and stay with the discomfort for five more minutes. Confusion is not the opposite of
understanding. It is the doorway to it.
\end{recoverybox}

\section*{A word on rigor, and a word of welcome}

You may have heard that calculus comes in two flavors: an easy ``intuitive'' version for engineers
and a hard ``rigorous'' version for mathematicians, and that you must choose between them. This book rejects that
choice. Intuition and rigor are not rivals; they are partners, and each is incomplete
without the other. Intuition without rigor led Newton and Leibniz to brilliant results along with
occasional nonsense, with no systematic way to distinguish them. Rigor without intuition is how calculus
becomes a dry sequence of formal $\eps$ and $\delta$ manipulations that produce correct proofs of theorems you cannot
picture and do not believe. We will develop both side by side. We use intuition to discover and pictures to
see, and we use rigor to be \emph{certain} -- to know, not merely to hope, that the geometric insight is
sound. The history boxes will show you, again and again, that this is how the subject
truly developed: the insight first, followed by the rigorous foundation.

\begin{historybox}
The people who invented calculus did not
understand it as well as you are about to. Newton and Leibniz, working independently in the 1660s and
1670s, had the central insights and the spectacular applications, but their foundations were genuinely
shaky -- they leaned on infinitesimal quantities nobody could define without contradiction. It took
roughly two more centuries for Augustin-Louis Cauchy and Karl Weierstrass to replace those ghosts
with the precise theory of limits you will learn in the early chapters. So if the foundations feel
subtle, that is not your weakness. They \emph{are} subtle -- subtle enough to elude the greatest minds
of the seventeenth and eighteenth centuries. You get to learn, in a few clear chapters, what cost
humanity two hundred years to work out. That is the gift of being born later.
\end{historybox}

\begin{connectionbox}
One last orientation. Although this book is about calculus, you will keep glimpsing the rest of
mathematics through its windows. The limit connects to the deep structure of the real numbers (an
appendix at the end pursues this). The derivative is a linear approximation, which ties calculus to
linear algebra. Series connect to the theory of approximation and to complex analysis. Differential
equations open onto all of physics. We will not chase every thread, but we will point at them, so you
always feel that calculus is one province of a single connected country, not an island.
\end{connectionbox}

So: welcome. You are about to learn one of the most powerful and beautiful things human beings have
ever figured out -- a way of thinking that decoded the planets, underlies every bridge and circuit
and economic model, and reveals a hidden unity beneath change and accumulation that nobody suspected
was there. It was hard-won, and it is worth winning. Read slowly. Draw constantly. Ask why without
mercy. Get stuck, and stay a little longer than is comfortable.

Let us begin from the ground up.

\section*{Exercises}

These are not computational drills -- the real exercises begin in Chapter~1. They are meant to fix
the habits and the central image of this preface before you turn the page.

\begin{enumerate}
  \item \textbf{The closing-in feeling.} Return to the ``Work it yourself'' box above, where the
  quotient $\dfrac{(1+h)^2-1}{h}$ simplified to $2+h$. Compute it once more, this time for
  $h=-0.1$, $h=-0.01$, and $h=-0.001$ -- that is, approaching from \emph{below} $x=1$. Do the values
  still close in on $2$? What does the simplified form $2+h$ predict they must do? (Approaching a
  value from both sides and getting the same answer is exactly what a limit will demand of us.)

  \item \textbf{Reading the picture.} Look carefully at Figure~\ref{ch00:secant-tangent}. Every gray
  secant passes through the same point $(1,1)$, yet each has a different slope. In one or two
  sentences, say in your own words what is sliding, what is staying fixed, and what the red line is
  the limit \emph{of}.

  \item \textbf{The arc.} Without looking back, try to redraw Figure~\ref{ch00:arc} from memory --
  the chain from functions all the way to differential equations. Which single box did the preface
  call the keystone, and why does every arrow after it secretly carry that same idea forward?

  \item \textbf{Carry the question.} Write down, in one sentence, the speedometer paradox from the
  ``Walking question'' box: how can a car have a definite speed at a single frozen instant, when no
  time passes and no distance is covered in an instant? You are not asked to solve it -- only to state
  it sharply. The whole first half of the book is the answer.
\end{enumerate}
